\SourceLeaf{076}
Pro aggregato \(S\) autem habemus formulam novam (20)
\[
S=\frac{\mathfrak A\mathfrak A}{[aa]}+\frac{\mathfrak B'\mathfrak B'}{[bb,1]}+\frac{\mathfrak C''\mathfrak C''}{[cc,2]}+\frac{\mathfrak D'''\mathfrak D'''}{[dd,3]}\quad\text{etc.}
\]
Denique si pondus \(P\), quod determinationi maxime plausibili quantitatis per functionem \(u\) expressae tribuendum est, desideratur, faciemus (21)
\[
\begin{aligned}\relax
[bl,1]&=[bl]-\frac{[ab][al]}{[aa]}\\[5pt]
[cl,2]&=[cl]-\frac{[ac][al]}{[aa]}-\frac{[bc,1][bl,1]}{[bb,1]}\\[5pt]
[dl,3]&=[dl]-\frac{[ad][al]}{[aa]}-\frac{[bd,1][bl,1]}{[bb,1]}-\frac{[cd,2][cl,2]}{[cc,2]}
\end{aligned}
\]
etc., quo facto erit (22)
\[
\frac1P=[ll]-\frac{[al]^2}{[aa]}-\frac{[bl,1]^2}{[bb,1]}-\frac{[cl,2]^2}{[cc,2]}-\frac{[dl,3]^2}{[dd,3]}-\text{etc.}
\]

Formulae (17)\(.\ .\ .\ .\)(22), quarum simplicitas nihil desiderandum relinquere videtur, solutionem problematis nostri ab omni parte completam exhibent.

\Article{14}
Postquam problemata primaria absolvimus, adhuc quasdam quaestiones secundarias attingemus, quae huic argumento maiorem lucem affundent.

Primo inquirendum est, num eliminatio. per quam \(x,y,z\) etc. ex \(\xi,\eta,\zeta\) etc. derivare oportet, umquam impossibilis fieri possit. Manifesto hoc eveniet, si functiones \(\xi,\eta,\zeta\) etc. inter se haud independentes essent. Supponamus itaque aliquantisper, unam earum per reliquas iam determinari, ita ut habeatur aequatio identica
\[
\alpha\xi+\beta\eta+\gamma\zeta+\text{etc.}=0
\]
denotantibus \(\alpha,\beta,\gamma\) etc. numeros determinatos. Erit itaque
\[
\begin{aligned}
\alpha[aa]+\beta[ab]+\gamma[ac]+\text{etc.}&=0\\
\alpha[ab]+\beta[bb]+\gamma[bc]+\text{etc.}&=0\\
\alpha[ac]+\beta[bc]+\gamma[cc]+\text{etc.}&=0
\end{aligned}
\]
etc., unde, si statuimus
\SourceLeaf{077}
\[
\begin{aligned}
\alpha a+\beta b+\gamma c+\text{etc.}&=p\theta\\
\alpha a'+\beta b'+\gamma c'+\text{etc.}&=p'\theta'\\
\alpha a''+\beta b''+\gamma c''+\text{etc.}&=p''\theta''
\end{aligned}
\]
etc., sponte sequitur
\[
\begin{aligned}
a\theta+a'\theta'+a''\theta''+\text{etc.}&=0\\
b\theta+b'\theta'+b''\theta''+\text{etc.}&=0\\
c\theta+c'\theta'+c''\theta''+\text{etc.}&=0
\end{aligned}
\]
etc., nec non
\[
p\theta\theta+p'\theta'\theta'+p''\theta''\theta''+\text{etc.}=0
\]
quae aequatio, quum omnes \(p,p',p''\) etc. natura sua sint quantitates positivae, manifesto consistere nequit, nisi fuerit \(\theta=0,\ \theta'=0,\ \theta''=0\) etc.

Iam consideremus valores differentialium completorum \(\mathrm dX,\mathrm dY,\mathrm dZ\) etc., respondentes valoribus iis quantitatum \(v,v',v''\) etc., ad quos referuntur observationes. Haec differentialia, puta
\[
\begin{aligned}
&a\mathrm dv+a'\mathrm dv'+a''\mathrm dv''+\text{etc.}\\
&b\mathrm dv+b'\mathrm dv'+b''\mathrm dv''+\text{etc.}\\
&c\mathrm dv+c'\mathrm dv'+c''\mathrm dv''+\text{etc.}
\end{aligned}
\]
etc., per conclusionem, ad quam modo delati sumus, inter se ita dependentia erunt, ut per \(\alpha,\beta,\gamma\) etc. resp. multiplicata aggregatum identice evanescens producant, sive quod idem est, quodvis ex ipsis (cui quidem respondet multiplicator \(\alpha,\beta,\gamma\) etc. non evanescens) sponte evanescet, simulac omnia reliqua evanescere supponuntur. Quamobrem ex aequationibus conditionalibus \(X=0,\ Y=0,\ Z=0\) etc., una (ad minimum) pro \textit{superflua} habenda est, quippe cui sponte satisfit, simulac reliquis satisfactum est.

Ceterum si res profundius inspicitur, apparet, hanc conclusionem per se tantum pro ambitu infinite parvo variabilitatis indeterminatarum valere. Scilicet proprie duo casus distinguendi erunt, alter, ubi una aequationum conditionalium \(X=0,\ Y=0,\ Z=0\) etc. absolute et generaliter iamiam in reliquis contenta est, quod facile in quovis casu averti poterit; alter, ubi, quasi fortuito, pro iis valoribus concretis quantitatum \(v,v',v''\) etc., ad quos observationes referrun\SourceLeaf{078}%
tur, una functionum \(X,Y,Z\) etc. e. g. prima \(X\), valorem maximum vel minimum (vel generalius, stationarium) nanciscitur respectu mutationum omnium, quas quantitatibus \(v,v',v''\) etc., salvis aequationibus \(Y=0,\ Z=0\) etc., applicare possemus. Attamen quum in disquisitione nostra variabilitas quantitatum tantummodo intra limites tam arctos consideretur, ut ad instar infinite parvae tractari possit, hic casus secundus (qui in praxi vix umquam occurret) eundem effectum habebit, quem primus, puta una aequationum conditionalium tamquam superflua reiicienda erit, certique esse possumus, si omnes aequationes conditionales retentae eo sensu, quem hic intelligimus, ab invicem independentes sint, eliminationem necessario fore possibilem. Ceterum disquisitionem uberiorem, qua hoc argumentum, propter theoreticam subtilitatem potius quam practicam utilitatem haud indignum est, ad aliam occasionem nobis reservare debemus.

\Article{15}
In commentatione priore art. 37 sqq. methodum docuimus, observationum praecisionem a posteriori quam proxime eruendi. Scilicet si valores approximati \(\pi\) quantitatum per observationes aequali praecisione gaudentes innotuerunt, et cum valoribus iis comparantur, qui e valoribus maxime plausibilibus \(\rho\) elementorum, a quibus illae pendent, per calculum prodeunt: differentiarum quadrata addere, aggregatumque per \(\pi-\rho\) dividere oportet, quo facto quotiens considerari poterit tamquam valor approximatus quadrati erroris medii tali observationum generi inhaerentis. Quoties observationes inaequali praecisione gaudent, haec praecepta eatenus tantum mutanda sunt, ut quadrata ante additionem per observationum pondera multiplicari debeant, errorque medius hoc modo prodiens ad observationes referatur, quarum pondus pro unitate acceptum est.

Iam in tractatione praesenti illud aggregatum manifesto quadrat cum aggregato \(S\), differentiaque \(\pi-\rho\) cum multitudine aequationum conditionalium \(\sigma\), quamobrem pro errore medio observationum, quarum pondus \(=1\), habebimus expressionem \(\BareRoot{\frac S\sigma}\), quae determinatio eo maiore fide digna erit, quo maior fuerit numerus \(\sigma\).

Sed operae pretium erit, hoc etiam independenter a disquisitione priore stabilire. Ad hunc finem quasdam novas denotationes introducere conveniet. Scilicet respondeant valoribus indeterminatarum \(\xi,\eta,\zeta\) etc. his
